% Удиртгал
% \phantomchapter{Удиртгал} % Зарим нэг зөвлөмж

\phantomsection
\addchaptertocentry{Удиртгал}
\label{Introduction} % Энэ бүлэг рүү ишлэл хийх бол \ref{Introduction} командыг ашигла 
%-------------------------------------------------------------------------------
%   INTRODUCTION
%-------------------------------------------------------------------------------
% Зорилго зорилт, өмнөх бие даалт харах
    \section*{Удиртгал}
    Мэдээлэл технологийн хурдацтай хөгжил нь боловсролын салбарт суралцах арга барилыг үндсээр нь өөрчилж байна. Уламжлалт тэмдэглэл хөтлөх арга нь үр дүн муутай, цаг их зарцуулдаг тул оюутнууд идэвхтэй суралцах арга барилд шилжих хэрэгцээ үүссэн. Хиймэл оюун ухааны том хэлний загварууд (LLMs) нь текстийг ойлгох, боловсруулах өндөр чадамжтай болсон нь энэхүү асуудлыг шийдвэрлэх технологийн суурь болж байна. Энэхүү ажил нь оюутны сургалтын материалыг интерактив, ухаалаг хэлбэрт шилжүүлэх системийг хөгжүүлэхэд чиглэнэ.

    \subsection*{Зорилго}
    Энэхүү төслийн зорилго нь их сургуулийн оюутнуудад зориулсан хиймэл оюунд суурилсан ухаалаг тэмдэглэлийн болон хувийн сургалтын менежментийн системийг хөгжүүлж, сургалтын материалыг автоматаар бүтэцжүүлэх, хураангуйлах болон суралцах ахицыг тоон үзүүлэлтээр хянах боломжийг бүрдүүлснээр суралцах үр ашгийг нэмэгдүүлэхэд оршино.
    
    \subsection*{Зорилт}
        Уг зорилгод хүрэхийн тулд дараах зорилтуудыг дэвшүүлж байна. Үүнд:
        \begin{enumerate}
            \item \textbf{Шаардлага тодорхойлох} \newline 
            Оюутан болон админ хэрэглэгчийн 2 үндсэн хэрэглэгчийн шаардлагад тулгуурлан системийн функциональ болон функциональ бус шаардлага, ижил төстэй системийн судалгааны баримт бичгийг боловсруулж, хэрэглэгчийн бүртгэл, нэвтрэлт бүхий хэрэглэгчийн удирдлагын модулийг хэрэгжүүлэх.
            \item \textbf{Сургалтын материалын автомат боловсруулалтын модулийг хэрэгжүүлэх} \newline 
            PDF боловсруулагч болон хиймэл оюуны моделын API ашиглан сургалтын материалыг хичээл → бүлэг → дэд сэдэв бүтэцтэйгээр автоматаар боловсруулж тэмдэглэл үүсгэн, тэмдэглэл бүр дээр хураангуй, мэдлэгийн карт болон тест  автоматаар үүсгэдэг модулийг хэрэгжүүлэх.
            \item \textbf{Хиймэл оюунт чатбот болон тестийн модулыг хэрэгжүүлэх} \newline 
            Хичээлийн агуулгад тулгуурласан асуулт-хариултын хиймэл оюунт чатбот-ыг хэрэгжүүлж, хэрэглэгч тухайн хичээлийн хүрээнд асуулт асууж, нэмэлт тайлбар авах болон тэмдэглэлд автоматаар үүссэн тест бүхий шалгалтын модулийг хэрэгжүүлэх.
            \item \textbf{Ахиц, аналитик, хяналтын самбар хэрэгжүүлэх} \newline 
            Хэрэглэгчийн суралцах явцыг бүртгэж, тестийн амжилт, хичээлийн төлөв (дуусгасан /үргэлжилж байгаа) зэрэг үзүүлэлтүүдийг харуулах хяналтын самбарыг 2026 оны 5-р сарын дотор хэрэгжүүлнэ.
            \item \textbf{Системийг байршуулж ашиглалтад оруулах (Deploy)} \newline 
            Хөгжүүлсэн веб системийг сервер орчинд байршуулж, бодит хэрэглэгч нэвтрэн ашиглах боломжтой болгож, үндсэн функцууд тогтвортой ажиллаж байгаа эсэхийг шалгах.
        \end{enumerate}

    \subsection*{Судлах үндэслэл}
        Сүүлийн үеийн судалгаагаар оюутнууд лекцийн материалыг зөвхөн унших (passive reading) нь мэдлэг тогтоох үр дүн багатай болох нь батлагдсан. Оюутнуудад их хэмжээний материалыг богино хугацаанд боловсруулах, өөрийгөө шалгах хэрэгсэл дутагдалтай байна. Одоо байгаа тэмдэглэлийн аппликэйшнүүд нь зөвхөн текст хадгалах үүрэгтэй бөгөөд агуулгыг шинжлэх, сургалтын хэрэглэгдэхүүн үүсгэх тал дээр хязгаарлагдмал байгаа нь энэхүү сэдвийг судлах үндэслэл болж байна.
        

    \subsection*{Судлагдсан байдал}
        Хиймэл оюун ухааныг боловсролд ашиглах чиглэлээр олон улсад өргөн хүрээний судалгаа хийгдэж байна. Ялангуяа автомат асуулт үүсгэх болон текстийг хураангуйлах технологиуд хөгжиж байгаа ч, эдгээрийг нэгдсэн нэг сургалтын платформ болгон оюутанд хүргэх шийдэл Монголын нөхцөлд хомс байна.
        
    \subsection*{Шинэлэг тал}
        Энэхүү системийн шинэлэг тал нь зөвхөн текст хураангуйлаад зогсохгүй, тухайн материалаас суралцахуйн шаталсан бүтэц үүсгэж, сэдэв тус бүрд тохирсон сорил болон флаш картыг автоматаар бэлтгэдэгт оршино. Мөн хэрэглэгчийн тестийн гүйцэтгэлээс хамаарч сэдвийг "Судалсан" гэж тэмдэглэх логик агуулсан нь бусад энгийн AI хэрэгслүүдээс ялгарна.

    \subsection*{Технологийн ач холбогдол}
        RAG (Retrieval-Augmented Generation) технологийг ашиглан хиймэл оюун ухааны "хий үзэгдэл" (hallucination)-ийг багасгаж, зөвхөн байршуулсан ном, лекцийн хүрээнд хариулт өгөх чатботыг хөгжүүлснээр технологийн найдвартай байдлыг харуулна.

        
    \subsection*{Хамрах хүрээ}
    Төсөл нь их, дээд сургуулийн бакалаврын түвшний оюутнуудад зориулагдсан бөгөөд Монгол болон Англи хэл дээрх PDF өргөтгөлтэй текстэн материалыг боловсруулах боломжтой вэб систем байна.